% AC-S49-ratio-simplicity: scoped necessity for the exact S48 ratio.
\section{The actual prolate ratio at a hypothetical multiple zero}
\label{sec:prolate-ratio-multiplicity}

This section tests the exact ratio (RZ.3), not a replacement power-law
lower bound.  The trial throughout is the original mean-zero prolate
combination (SR.2); no truncated copy of its limiting theta kernel is
substituted for it.  The conclusion is only a necessary consequence of
that ratio.  It supplies neither a second-eigenvalue lower bound nor a
proof of the ratio from \(A\).

Put \(c=2\pi\lambda^2\), \(L=\log\lambda\), and
\(\delta_j=1-\nu_j\).  For a fixed even prolate index \(j\), let
\[
 Y_j(x)=\sqrt\lambda\,(\mathcal F\phi_{j,\lambda})(\lambda x),
 \quad E_j=Y_j(1),\quad x\ge1.
 \tag{PM.1}
\]
The exterior differential equation, including its regular endpoint,
is the one proved in AC-S31:
\[
 ((x^2-1)Y_j')'+(c^2x^2-\chi_j)Y_j=0,
 \qquad 0\le\chi_j\le C_jc,
 \qquad\int_1^\infty |Y_j|^2dx=\delta_j/2.
 \tag{PM.2}
\]
For clarity every estimate below is for fixed indices, here only
\(j=0,4\).  No growing-index prolate assertion is used.

\begin{theorem}[AC-S49-ratio-simplicity: the exact ratio implies simplicity]
\label{AC-S49-ratio-simplicity}
Assume \(A\).  If the actual original trial satisfies
\[
 \frac{q_L(v_\lambda)}{a_2(L)}\longrightarrow0,
 \tag{PM.3}
\]
where \(v_\lambda\) is the normalized real even trial of AC-S48,
then every nontrivial zeta zero is simple.

More precisely, if \(\gamma>0\) were a multiple zero, there would
exist an unbounded sequence of actual parameters \(\lambda\) on which
\[
 |\widehat v_\lambda'(\gamma)|^2
       \le C_\gamma L^2 q_L(v_\lambda).
 \tag{PM.4}
\]
Along that sequence the ratio in (PM.3) is bounded below by a positive
constant.  This is a conditional obstruction; no multiple zero or
logical independence from \(A\) is asserted.
\end{theorem}

\begin{proof}
\proofstep{1. A two-sided endpoint scale and a global derivative bound.}
We first sharpen the exterior upper estimate of AC-S31 by using both
sides of the same oscillator-energy comparison.  Write \(E=Y_j(1)\).
It is nonzero: the regular endpoint equation would otherwise give
\(Y_j'(1)=0\), and its regular solution would vanish identically.
The energy argument of (SE.4)--(SE.5), now in dimensionless coordinates,
gives
\[
 |Y_j(x)|\le |E|,\qquad |Y_j'(x)|\le Cc^2|E|
                       \quad(1\le x\le2).
 \tag{PM.5}
\]
Set \(x=\cosh t\), \(Z(t)=\sqrt{\sinh t}\,Y_j(\cosh t)\),
and \(t_0=a/c\), where \(a>0\) is fixed and sufficiently small.
Exactly as in (SE.6),
\[
 Z''+\Omega Z=0,\qquad
 \Omega=c^2\cosh^2t-\chi_j-\tfrac14+\frac1{4\sinh^2t}.
 \tag{PM.6}
\]
On \([t_0,t_1]\), \(t_1=\operatorname{arcosh}2\), \(\Omega\)
is between positive constant multiples of \(c^2\), and first
decreases and then increases.  At \(t_0\), (PM.5) gives
\[
 |Z(t_0)|\le C_a c^{-1/2}|E|,\quad
 |Z'(t_0)|\le C_a c^{1/2}|E|,\quad
 |Z'(t_0)|^2+\Omega(t_0)|Z(t_0)|^2\le C_ac|E|^2.
\]
For \(\mathcal E_Z=|Z'|^2+\Omega|Z|^2\),
\(\mathcal E_Z'=\Omega'|Z|^2\).  On a decreasing portion the
energy does not increase, and on an increasing portion
\(\mathcal E_Z'\le(\Omega'/\Omega)\mathcal E_Z\).
The bounded ratio of the extrema of \(\Omega\) therefore implies
\[
 \mathcal E_Z(t)\le Cc|E|^2,
 \qquad |Y_j(\cosh t)|^2\le\frac{C|E|^2}{c\sinh t}
                       \quad(t_0\le t\le t_1).
 \tag{PM.7}
\]
For \(0<t<t_0\) use (PM.5).  Integrating with \(dx=\sinh t\,dt\)
now gives \(\int_1^2|Y_j|^2dx\le C|E|^2/c\).

For \(x\ge2\), set \(U=xY_j\).  Its equation is
\[
 U''+bU'+k_c^2U=0,\qquad b=\frac2{x(x^2-1)},\quad
 k_c^2=c^2+\frac{c^2-\chi_j-2/x^2}{x^2-1}.
 \tag{PM.8}
\]
Here \(k_c^2\ge c^2\) and \((k_c^2)'<0\) for large \(c\),
so \(|U'|^2+k_c^2|U|^2\) decreases.  At \(x=2\) this energy
is comparable with the energy in (PM.7), with absolute fixed-index
constants.  Thus
\[
 |Y_j(x)|\le\frac{C|E_j|}{\sqrt c\,x},\qquad
 |Y_j'(x)|\le\frac{C\sqrt c\,|E_j|}{x}
                      \quad(x\ge2).
 \tag{PM.9}
\]
Integrating the first estimate and combining the preceding interior
bound gives \(\delta_j\le C|E_j|^2/c\).  The reverse comparison
\(|E_j|^2\le Cc\delta_j\) is already (SE.11), after scaling.
Consequently
\[
 |E_j|\asymp\sqrt{c\delta_j},\qquad
 |Y_j(x)|\le C\sqrt{\delta_j}/x,\quad
 |Y_j'(x)|\le Cc\sqrt{\delta_j}/x\quad(x\ge2).
 \tag{PM.10}
\]
No derivative of the Fuchs eigenvalue asymptotic enters this step.

\proofstep{2. Fixed exterior phase from a regular Bessel matching.}
The Hermite comparison already used in AC-S30 implies
\[
                       \chi_j/c\longrightarrow2j+1.
 \tag{PM.11}
\]
Here is a direct verification without a derivative estimate for the
Hermite approximation.  Divide the time-interval prolate equation by
\(\lambda^2\), test it against a fixed compact smooth function, and
integrate its differential operator onto that test.  The unit-norm
prolates converge in \(L^2\) to their Hermite functions, and
\(\chi_j/\lambda^2\) is bounded by (PM.2).  The limiting operator is
\(-\partial_u^2+4\pi^2u^2\), whose \(j\)-th Hermite eigenvalue is
\(2\pi(2j+1)\).  A test with nonzero Hermite pairing proves (PM.11).

Normalize \(Y_j\) by its endpoint \(E_j\), put
\(z=\sqrt{x^2-1}\), and write
\[
 Y_j(x)/E_j=x^{-1/2}v(z),\qquad w(z)=\sqrt z\,v(z).
\]
Direct differentiation gives
\[
 v''+z^{-1}v'+(c^2-q_c(z))v=0,\quad v(0)=1,\ v'(0)=0,
 \quad q_c(z)=\frac{\chi_j}{1+z^2}
                  +\frac{z^2+4}{4(1+z^2)^2}.
 \tag{PM.12}
\]
Equivalently \(w''+(c^2-q_c+1/(4z^2))w=0\).
We give the matching error because an unproved exterior asymptotic
would not suffice for the present subsequence argument.

On \(0\le z\le1/c\), variation of constants for the regular Bessel
solution is
\[
 v(z)=J_0(cz)+\frac\pi2\int_0^z t
       [J_0(ct)Y_0(cz)-Y_0(ct)J_0(cz)]q_c(t)v(t)\,dt.
 \tag{PM.13}
\]
The Bessel kernel is bounded there by
\(Ct(1+|\log(z/t)|)\); its \(z\)-derivative is bounded by
\(C(t/z)(1+|\log(z/t)|)\).
These follow directly from the convergent small-argument series.
Since \(q_c=O(c)\), the integral operator in (PM.13) has sup norm
\(O(cz^2)\), at most \(O(c^{-1})\).  Thus \(v\) is bounded on
this interval, and at \(z=1/c\) its value and its derivative divided
by \(c\) differ from those of \(J_0(cz)\) by \(O(c^{-1})\).

On \([1/c,z_0]\), \(z_0=c^{-1/2}\), use the fundamental matrix
with columns
\[
 \begin{pmatrix}\sqrt z J_0(cz)\\(\sqrt z J_0(cz))'/c\end{pmatrix},
 \qquad
 \begin{pmatrix}\sqrt z Y_0(cz)\\(\sqrt z Y_0(cz))'/c\end{pmatrix}.
\]
Its norm is \(O(c^{-1/2})\), and its inverse has norm
\(O(c^{1/2})\); the determinant is exactly \(2/(\pi c)\).
Variation of constants for the added potential \(q_c\) therefore
has coefficient-matrix norm \(O(q_c/c)=O(1)\).  Gronwall on an
interval of length at most \(z_0\), together with the initial error,
shows that at \(z_0\),
\[
 \begin{pmatrix}w\\w'/c\end{pmatrix}
 =\sqrt{\frac2{\pi c}}
       \begin{pmatrix}\cos(cz_0-\pi/4)\\-\sin(cz_0-\pi/4)\end{pmatrix}
       +O(c^{-1}).
 \tag{PM.14}
\]
The error in the usual fixed-order large-argument Bessel formula is
relative \(O((cz_0)^{-1})=O(c^{-1/2})\), included here.

On \([z_0,R]\), for fixed \(R>0\), put
\(k(z)=(c^2-q_c(z)+1/(4z^2))^{1/2}\).
It is positive and comparable with \(c\), and direct integration
of its derivative gives
\[
                  \int_{z_0}^R |k'(z)|/k(z)\,dz=O(c^{-1}).
 \tag{PM.15}
\]
Indeed \(q_c'=O(cz/(1+z^2)^2)+O(1)\), while the extra derivative
is \(O(z^{-3})\); division by \(c^2\) and integration give the
stated bound.  For the vector \((\sqrt k\,w,w'/\sqrt k)\), rotation
by \(\int k\) removes the skew part of its first-order equation.
The remaining matrix has norm \(|k'|/(2k)\), so (PM.15) controls
its full accumulated error.  Moreover Taylor expansion of the square
root, with its integrated remainder, gives
\[
 \int_{z_0}^z k(t)\,dt
  =c(z-z_0)-\frac{\chi_j}{2c}(\arctan z-\arctan z_0)+o(1).
 \tag{PM.16}
\]
Uniformly for \(z\) in any fixed positive compact interval, the
omitted singular phase is \(O((cz_0)^{-1})\), and the integrated
quadratic error is
\(O(c^{-1}+c^{-3}z_0^{-3}+c^{-2}z_0^{-1})=o(1)\).
Equations (PM.11), (PM.14)--(PM.16), followed by the exact change
back to \(x\), prove, uniformly on fixed compact subsets of \((1,\infty)\),
\[
 \frac{Y_j'(x)}{E_j\sqrt{2c/\pi}}
 =-\frac{x^{1/2}}{(x^2-1)^{3/4}}
       \sin\left(c\sqrt{x^2-1}
           -\frac{2j+1}{2}\arctan\sqrt{x^2-1}-\frac\pi4\right)+o(1).
 \tag{PM.17}
\]
In particular this is an actual fixed-mode exterior statement, not an
analytic continuation of a Hermite error bound.

The only classical Bessel inputs in this proof are its two elementary
series, Wronskian, and fixed-order first asymptotic terms (also for the
first derivative).  Exact locators are given after the proof.  The
prolate matching, its errors and its parameter range were proved here.

\proofstep{3. Exact Mellin expansion of the actual lower theta tail.}
For \(s\) in a fixed compact subset of \(0<\operatorname{Re}s<1\), define
\[
 H_s(x)=\sum_{1\le n\le x}n^{s-1},\qquad
 f_s(x)=x^{-s}H_s(x)-1/s\quad(x\ge1),\qquad p(x)=x^2-1.
 \tag{PM.18}
\]
The values at integer endpoints do not affect its integrals.  Comparing
the sum with its integral gives, uniformly on such compact sets,
\(f_s(x)=O(x^{-a})\) for some \(a>0\), with an extra fixed power
of \(1+\log x\) after differentiating in \(s\).  On open unit
intervals and at integers \(n\ge2\), respectively,
\[
 f_s'=-s f_s/x-1/x,\qquad [f_s]_n=1/n.
 \tag{PM.19}
\]
Here brackets are right minus left limits.
For \(k_s=f_s/x^2\),
\[
 [k_s]_n=n^{-3},\qquad [k_s']_n=-(s+2)n^{-4},
\]
and a direct calculation on the open intervals gives
\[
 (pk_s')'=
 (s+2)(s+1)f_sx^{-2}-(s+2)(s+3)f_sx^{-4}
             +(s+3)x^{-2}-(s+5)x^{-4}.
 \tag{PM.20}
\]
The functions on this right side and their first \(s\)-derivatives
are uniformly in \(L^2(1,\infty)\).  The same is true with \(f_s/x^4\)
in place of \(k_s\).

Apply integration by parts separately on the unit intervals to (PM.2).
The boundary at \(1\) vanishes because \(p(1)=0\); the boundary at
infinity vanishes by (PM.9) and the decay of \(k_s,k_s'\).  Summing
the interior boundary terms gives the exact identity
\[
\begin{split}
 c^2\int_1^\infty Y_j f_s\,dx
 &=\chi_j\int_1^\infty Y_j f_sx^{-2}\,dx
     -\int_1^\infty Y_j(pk_s')'_{\rm pw}\,dx\\
 &\quad+\sum_{n\ge2}p(n)
       \{n^{-3}Y_j'(n)+(s+2)n^{-4}Y_j(n)\}.
\end{split}
 \tag{PM.21}
\]
The subscript denotes the ordinary derivative on each open interval;
all jump terms are already displayed.  They converge absolutely by
(PM.10), as do their first \(s\)-derivatives.  Applying exactly the
same identity to \(f_s/x^2\), and using its jumps \(n^{-3}\) before
the second division by \(x^2\), gives
\[
 \left|\int_1^\infty Y_j f_sx^{-2}\,dx\right|
  +\left|\partial_s\int_1^\infty Y_j f_sx^{-2}\,dx\right|
                         \le C\sqrt{\delta_j}/c.
 \tag{PM.22}
\]
To see the scale explicitly, its undivided bulk terms are at most
\(C(1+\chi_j)\sqrt{\delta_j}\), and its derivative jumps are bounded
by \(Cc\sqrt{\delta_j}\sum_{n\ge2}n^{-4}\); division by \(c^2\)
gives (PM.22).  Thus (PM.21) implies
\[
 \int_1^\infty Y_j f_s\,dx=A_j(c)+O(\sqrt{\delta_j}/c^2),\qquad
 \partial_s\int_1^\infty Y_j f_s\,dx=O(\sqrt{\delta_j}/c^2),
 \quad
 A_j(c)=c^{-2}\sum_{n\ge2}\frac{n^2-1}{n^3}Y_j'(n).
 \tag{PM.23}
\]
The leading term is real and independent of \(s\), which is the key
point.  All error constants are uniform on the specified \(s\)-compact.

We now identify (PM.23) with the actual theta tail, including its scalar.
Write \(\psi_\lambda=\alpha_0\phi_{0,\lambda}+\alpha_4\phi_{4,\lambda}\)
with the exact coefficients in (SR.2), and set
\[
 Y=\alpha_0Y_0+\alpha_4Y_4,\qquad
 R(t)=E\psi_\lambda(e^{-L-t}),\quad t>0.
\]
The mean of \(\psi_\lambda\) is exactly zero.  For \(x=e^t>1\),
Poisson summation in its Fejer sense gives
\[
 R(t)=e^{t/2}\sum_{n\ge1}Y(ne^t)
              -\frac{\psi_\lambda(0)}{2\sqrt\lambda}e^{-t/2}.
 \tag{PM.24}
\]
No absolute convergence of this theta series is asserted.  Here is a
precise integral justification.  The one-period function
\[
 P_x(v)=\sum_{m\in\mathbb Z}
                 \psi_\lambda((m+v)/(\lambda x))
\]
has Fourier coefficients \(\lambda x\,\mathcal F\psi_\lambda(n\lambda x)\).
Its mean is zero.  The elementary shifted-Riemann-sum bound gives
\(\|P_x\|_\infty\le\operatorname{Var}(\psi_\lambda)\), independently
of \(x\).  Its Fejer sums have the same bound, because the Fejer kernel
is nonnegative of integral one.  At every \(x\) except a countable set
of endpoint discontinuities their values at zero converge to \(P_x(0)\).
Since the function and its Fourier coefficients are even, its nonzero
positive Fourier terms are half that Fejer sum.  After the factor
\(e^{t/2}/(\sqrt\lambda e^t)\), they are dominated by
\(C_\lambda e^{-t/2}\).  This is integrable with every fixed first
logarithmic moment and with the compact \(s\)-weights below.
Dominated convergence therefore justifies the Mellin integral of
(PM.24) as a limit of finite Fejer sums.

For each mode, one further integration by parts with the test
\(x^{-s-2}\) on \([n,\infty)\) and (PM.10) gives
\[
 \left|n^{s-1}\int_n^\infty Y_j(x)x^{-s}\,dx\right|
                 \le C\sqrt{\delta_j}/(c n^2)\quad(n\ge2),
 \tag{PM.25}
\]
with summable logarithmic factors after one \(s\)-derivative.
The \(n=1\) integral converges too.  Thus the Fejer-weighted series
of these integrals converges absolutely to the unweighted one.
Writing its partial sums as an integral up to \(n\) and the remaining
tail shows that this is
\(\int_1^\infty Yf_s+(1/s)\int_1^\infty Y\).
For example the leftover term is bounded by \(C_\lambda/n\) using
(PM.25), so no conditional interchange is hidden here.
The latter ordinary improper integral exists by the same integration
by parts at \(s=0\).  Fourier inversion at the smooth raw origin gives
\[
 \int_1^\infty Y(x)\,dx
   =\lambda^{-1/2}\left(\frac{\psi_\lambda(0)}2-c_Q\right),
 \qquad c_Q=\int_0^\lambda\mathcal F\psi_\lambda(u)\,du.
\]
It follows, with the exact \(c_Q\) of (SR.18), that
\[
 \mathcal R_\lambda(s):=\int_0^\infty
           R(t)e^{-(s-1/2)t}\,dt
 =\sum_{j=0,4}\alpha_j\int_1^\infty Y_jf_s\,dx
                       -\frac{c_Q}{\sqrt\lambda\,s}.
 \tag{PM.26}
\]
The moment identity and its first \(s\)-derivative are consequently
valid for the actual lower tail.

The coefficients \(\alpha_j\) are bounded, \(\alpha_4\) is bounded
away from zero, and \(|c_Q|\le C\delta_4\), by AC-S30.  Its fixed-index
leakage input also gives \(\delta_0/\delta_4=O(c^{-4})\).
Equations (PM.23) and (PM.26) thus imply, for fixed compact \(s\)-sets,
\[
 \mathcal R_\lambda(s)=A(c)+O(\sqrt{\delta_4}/c^2),\qquad
 \partial_s\mathcal R_\lambda(s)=O(\sqrt{\delta_4}/c^2),\qquad
 A(c)=\alpha_0A_0(c)+\alpha_4A_4(c)\in\mathbb R.
 \tag{PM.27}
\]
The scalar term is absorbed here because
\(\delta_4\lambda^{-1/2}=o(\sqrt{\delta_4}/c^2)\), by the same
already-declared exponential leakage asymptotic.  This does not alter
or omit the polar correction.

\proofstep{4. A nonzero oscillatory leading coefficient.}
Put \(z_n=\sqrt{n^2-1}\) and
\[
 S(c)=\sum_{n\ge2}\frac{(n^2-1)^{1/4}}{n^{5/2}}
       \sin\left(cz_n-\frac92\arctan z_n-\frac\pi4\right),
 \qquad
 B(c)=-\alpha_4 E_4 c^{-3/2}\sqrt{2/\pi}.
 \tag{PM.28}
\]
This sine series converges absolutely and uniformly, since its
coefficients are \(O(n^{-2})\).  Equations (PM.17) and (PM.23)
give, first for each finite partial sum and then for the full sum,
\[
                  A(c)=B(c)(S(c)+o(1)),
 \qquad |B(c)|\asymp\sqrt{\delta_4}/c.
 \tag{PM.29}
\]
The passage to the infinite sum is uniform: (PM.9) bounds its normalized
tail by \(C\sum_{n>N}n^{-2}\), independently of \(c\).  The mode-zero
contribution divided by \(B(c)\) is \(O(c^{-2})\), using (PM.10) and
\(\delta_0/\delta_4=O(c^{-4})\).  Therefore neither an unproved
infinite-index asymptotic nor a pointwise nonzero phase assumption is
needed.

The continuous function \(S\) has positive mean square uniformly on
long intervals.  For a fixed finite partial sum, integrating its square
shows that its average on \([C,C+H]\) tends, uniformly in the start
\(C\), to half the sum of the squared coefficients as \(H\to\infty\).
All frequencies \(z_n\) are positive and distinct, so every cross term
has nonzero frequency and average \(O_N(H^{-1})\).  Uniform convergence
of the series transfers this assertion to a fixed positive lower bound
for the mean square of \(S\) on every sufficiently long interval.
In particular there is a fixed \(s_0>0\) such that every sufficiently
long interval contains a point where \(|S(c)|\ge s_0\).

For a fixed \(\gamma>0\), take the unbounded intervals defined by
\[
 |\gamma L-2\pi m|\le\pi/3,\qquad c=2\pi e^{2L}.
 \tag{PM.30}
\]
Their lengths in the \(c\)-variable tend to infinity, and
\(|\cos(\gamma L)|\ge1/2\) throughout each interval.  Choose in each
sufficiently late interval a point where \(|S(c)|\ge s_0\).
Equations (PM.27)--(PM.29) then hold at those actual parameters with
\(|A(c)|\ge(s_0/2)|B(c)|\).  They also give the global upper bounds
\(|\mathcal R_\lambda(s)|\le C|B(c)|\) and
\(|\partial_s\mathcal R_\lambda(s)|=o(|B(c)|)\) on the fixed compact
\(s\)-sets.  No Diophantine approximation of zeta ordinates is used.

\proofstep{5. Values, derivatives and the second-eigenvalue obstruction.}
Assume now that \(\gamma\) is a multiple zero.  The full uncut theta
function has value and first Fourier derivative zero there, because
its Mellin transform has the factor \(\zeta\); this is the exact
jet fact used in AC-S34.  Its upper tail is zero.  Therefore, with
\(N_\lambda=\|k_\lambda^+\|_2\to\|K\|_2>0\),
\[
 \widehat v_\lambda(\gamma)
       =-N_\lambda^{-1}\operatorname{Re}
           \{e^{i\gamma L}\mathcal R_\lambda(1/2-i\gamma)\}.
 \tag{PM.31}
\]
Differentiating before evaluation uses the vanishing full first jet
and gives the same formula with the derivative of the bracket.
By (PM.27), on the sequence selected in (PM.30),
\[
 |\widehat v_\lambda(\gamma)|\ge c_\gamma |B(c)|,
 \qquad
 |\widehat v_\lambda'(\gamma)|\le C_\gamma L|B(c)|.
 \tag{PM.32}
\]
The sampling form contains the value at \(\gamma\) with its positive
multiplicity.  Hence \(q_L(v_\lambda)\ge|\widehat v_\lambda(\gamma)|^2\),
which proves (PM.4).

For completeness, the short variational implication from (PM.4) is
included, rather than assumed as a gap estimate.  Fix real even
\(b\in C_c^\infty(-h,h)\) with \(\widehat b(\gamma)\ne0\), and
\(a_0>h+1\), \(a_1=a_0+\pi/(2\gamma)\).
The two even translates
\(b(x-L+a_i)+b(x+L-a_i)\) lie in the window for large \(L\).
Their value/derivative matrix at \(\gamma\), after dividing the second
row by \(L\), has uniformly bounded inverse: its original determinant
is exactly
\[
 4\widehat b(\gamma)^2
   \{L-a_1\cos^2((L-a_0)\gamma)-a_0\sin^2((L-a_0)\gamma)\}.
\]
Their form norms are bounded independently of \(L\), since translation
is a phase on the real zero samples.  Correcting the two jets thus
produces an even \(j_\lambda\) with exact double zeros at \(\pm\gamma\)
and
\[
 q_L(j_\lambda)\le C_\gamma
     \left(q_L(v_\lambda)+L^{-2}|\widehat v_\lambda'(\gamma)|^2\right).
 \tag{PM.33}
\]
AC-S47's absolute jet bound makes this an exponentially small \(L^2\)
correction, so \(j_\lambda\to K/\|K\|_2\) at rate
\(O(\lambda^{-1/2})\).  The compact sine inverse
\[
 g_\lambda(x)=-\gamma^{-1}\int_{-L}^x
           \sin(\gamma(x-u))j_\lambda(u)\,du\quad(|x|\le L),
\]
extended by zero, has transform
\(\widehat j_\lambda(t)/(t^2-\gamma^2)\).
Its boundary value and first derivative are zero, its transform still
vanishes at the removed ordinates, and it is in the closed form domain,
exactly as in (SM.8)--(SM.12).  The quantitative inverse bound
\(2L/\gamma\) and the rapid tails of \(K\) show convergence to the
inverse transform of \(\Xi/(\|K\|_2(t^2-\gamma^2))\).
This limit and \(K/\|K\|_2\) are independent, so their Gram matrices
are bounded below.  Discreteness at the single fixed ordinate bounds
all remaining division denominators below.  Positive-form
Cauchy--Schwarz and the two-dimensional minimum principle yield
\[
 a_2(L)\le C_\gamma
       \left(q_L(v_\lambda)+L^{-2}|\widehat v_\lambda'(\gamma)|^2\right).
 \tag{PM.34}
\]
Combining (PM.4) and (PM.34) on the chosen unbounded sequence bounds
\(a_2(L)/q_L(v_\lambda)\) above.  This contradicts (PM.3), proving
the asserted necessary simplicity condition.
\end{proof}

\paragraph{Exact source scope.}
The original trial, its coefficients, mean-zero identity and scalar
\(c_Q\) are exactly AC-S30, (SR.1)--(SR.2) and (SR.18).
The exterior equation and the monotone energies come from the complete
proof of AC-S31, freshly reread.  The two-sided upper passage, regular
Bessel matching, jump expansion, Fejer interchange, oscillatory
subsequence and relative comparison above are new calculations here.
The only fixed-index leakage asymptotic is the one already recorded in
AC-S30 with its original-proof OCR limitation; no differentiated version
is assumed.  AC-S34's compact inverse and AC-S47's jet estimates are
used at their complete proved scopes.

For the elementary Bessel inputs, the directly checked NIST DLMF
locators are (10.2.2), the first-kind series; (10.5.2), the Wronskian;
(10.6.3), first derivatives;
(10.7.1), small-argument behavior; (10.8.1)--(10.8.2), the convergent
series giving the kernel bounds; and (10.17.3)--(10.17.4) and (10.17.9)--(10.17.10), the
fixed-order large-positive-argument expansions and their derivatives.  These displayed
published formulas are retained inputs.  A fresh read of Olver's or
Watson's whole upstream proof is not claimed.  The proof above supplies
its own prolate perturbation and all its error bounds.

\paragraph{Scope of the result.}
The implication (PM.3) implies simple zeros is a necessary obstruction
to deriving the exact ratio from \(A\) alone by the current argument.
It proves neither that \(A\) fails to imply simplicity nor that the
specific ground-line conjecture is false.  No generic countermodel is
substituted for the actual zeta/prolate trial, and no new motivic or
etale realization is asserted.  The central positive task---a valid
lower-gap estimate or another actual proof of ground-line selection---
remains open.


\paragraph{Grothendieck-language dictionary.}
The previously audited finite-cover dictionary is SGA~1,
Expos\'e~V, \S\S4 and~7: finite etale covers and their fiber functors
supply the categorical language for finite Galois actions and finite
local systems.  We reuse that locator only with its established scope.
The present argument concerns a real Hilbert-space quadratic form,
an actual prolate differential equation, and analytic division by the
factor \(t^2-\gamma^2\).  That division and the positive-form minimum
principle do not construct a finite etale cover or a motivic quotient.
No new Grothendieck realization, arithmetic square, or cohomological
positivity theorem is claimed; the necessary consequence is precisely
the analytic simplicity implication in (PM.3)--(PM.4).
